\subsection{无单位元代数的情形}

定义 6. --- 设 A 是 K 上的一个交换代数。A 的特征标是指从 A 到 K 的一个代数同态。

A 的特征标所成的集合将记作 X'(A)。

零映射是一个代数同态。若 A 有单位元 \(e\)，则从 A 到 K 中的每个非零代数态射都是保单位元的，也就是说，它是定义 4 意义下的 A 的一个幺特征标：事实上，对 \(\chi \in \mathsf{X}'(\mathsf{A})\)，\(\chi\) 非零必须且只需 \(\chi(e) = 1\)。

我们令 \(\mathsf{X}(\mathsf{A})=\mathsf{X}^{\prime}(\mathsf{A})-\{0\}\)；根据上文，这一记号与 A 为幺代数时引入的记号相容。

若 \(h\colon A\to B\) 是交换代数的一个态射，则映射 \(\chi\mapsto\chi\circ h\) 是一个映射 \(\mathsf{X}^{\prime}(h)\colon\mathsf{X}^{\prime}(\mathrm{B})\to\mathsf{X}^{\prime}(\mathrm{A})\)。它把 0 映到 0。若 \(k\colon B\to C\) 是交换代数的一个态射，则有 \(\mathsf{X}^{\prime}(k\circ h)=\mathsf{X}^{\prime}(h)\circ\mathsf{X}^{\prime}(k)\)。若 \(h\) 是满射，则 \(\mathsf{X}^{\prime}(h)\) 是从 \(\mathsf{X}^{\prime}(\mathrm{B})\) 到在 \(h\) 的核上为零的 A 的特征标所成集合上的一个双射。设 \(\mathrm{A}_{1},\ldots,\mathrm{A}_{n}\) 是交换代数，\(\mathrm{A}=\mathrm{A}_{1}\times\cdots\times\mathrm{A}_{n}\)，而 \(\pi_{i}\colon\mathrm{A}\to\mathrm{A}_{i}\) 是典范态射；则 \(\mathsf{X}^{\prime}(\pi_{i})\) 是从 \(\mathsf{X}^{\prime}(\mathrm{A}_{i})\) 到子集 \(\mathsf{X}_{i}^{\prime}\)（\(\mathsf{X}^{\prime}(\mathrm{A})\) 的一个子集）上的双射，这个子集即由 A 的在 \(\prod_{j\neq i}\mathrm{A}_{j}\) 上为零的特征标所成的集合；我们如第 6 小节那样看出，\(\mathsf{X}^{\prime}(\mathrm{A})\) 是

诸 \(X_{i}^{\prime}\) 的并；另一方面，有 \(X_{i}^{\prime} \cap X_{j}^{\prime} = \{0\}\)（对 \(i \neq j\)）；特别地，诸 \(X_{i}^{\prime} - \{0\}\) 构成 \(X^{\prime}(A) - \{0\} = X(A)\) 的一个划分。

对每个 \(x \in A\)，令 \(\mathcal{G}_{\mathrm{A}}^{\prime}(x)\)（或简作 \(\mathcal{G}^{\prime}(x)\)）为映射 \(\chi \mapsto \chi(x)\)（从 \(\mathsf{X}^{\prime}(\mathrm{A})\) 到 K）。映射 \(\mathcal{G}^{\prime}\) 是从 A 到代数 \(A_{1}\)（由在 0 处为零的映射 \(\mathsf{X}^{\prime}(\mathrm{A}) \to \mathrm{K}\) 所成的代数）中的一个态射。设 B 是一个交换代数，\(B_{1}\) 是由在 0 处为零的映射 \(\mathsf{X}^{\prime}(\mathrm{B}) \to \mathrm{K}\) 所成的代数，h 是从 A 到 B 的一个态射；则 \(\mathsf{X}^{\prime}(h)\) 定义了一个态射 \(h_{1}: A_{1} \to B_{1}\)，且我们有 \(h_{1} \circ \mathcal{G}_{A}^{\prime} = \mathcal{G}_{B}^{\prime} \circ h\)。用 \(\mathcal{G}_{\mathrm{A}}(x)\)（或简作 \(\mathcal{G}(x)\)）表示 \(\mathcal{G}_{\mathrm{A}}^{\prime}(x)\) 在 \(\mathrm{X}(\mathrm{A})\) 上的限制，并称它为 \(x\) 的 Gelfand 变换。

设 \(\widetilde{\mathsf{A}}\) 是由 A 添加单位元而得到的幺代数。通过限制，\(\widetilde{\mathsf{A}}\) 的每个特征标都定义 A 的一个特征标；反之，A 的每个特征标可唯一地延拓为 \(\widetilde{\mathsf{A}}\) 的特征标。这就定义了一个从 \(\mathsf{X}'(\mathsf{A})\) 到 \(\mathsf{X}(\widetilde{\mathsf{A}})\) 上的典范双射，据此把这两个集合视为同一。A 的特征标 0 与 \(\widetilde{\mathsf{A}}\) 的以 A 为核的唯一特征标视为同一。

若 \(x\in \mathrm{A}\) 且 \(\chi \in \mathsf{X}'(\mathrm{A})\)，则有 \(\chi (x)\in \mathrm{Sp}_{\widetilde{\mathrm{A}}}^{\sim}(x)\)，从而 \(\chi (x)\in \mathrm{Sp}_{\mathrm{A}}^{\prime}(x)\)。

引理 3. --- 映射 \(\chi \mapsto \mathrm{Ker}(\chi)\) 是从 \(\mathsf{X}(\mathsf{A})\) 到 A 的余维数为 1 的正则理想所成的集合上的一个双射。

回顾 (A, VIII, p.~426, 定义 1)，A 的理想 I 称为正则的，如果商代数 A/I 有单位元。

我们来证明该引理。一方面，\(\mathsf{X}(\mathsf{A})\) 与 \(\widetilde{A}\) 的在 A 上非零的特征标所成的集合视为同一。另一方面，由 A, VIII, p.~428, 命题 4，映射 \(I \mapsto A \cap I\) 是从 \(\widetilde{A}\) 的异于 A 的极大理想所成的集合到 A 的正则极大理想所成的集合上的一个双射。于是引理由第 6 小节的结果得出。

现在设 K 是一个拓扑域。这时我们给 \(\mathsf{X}^{\prime}(\mathsf{A})\) 赋予 A 上逐点收敛的拓扑；记号 \(\mathsf{X}^{\prime}(\mathsf{A})\) 此后表示这样得到的拓扑空间。当 K = R 或 C 时，我们也称它为弱拓扑。对每个 \(x \in A\)，函数 \(\mathcal{G}_{\mathrm{A}}^{\prime}(x)\)（定义在 \(\mathsf{X}^{\prime}(\mathsf{A})\) 上）连续。

若 h 是从 A 到 B 的一个态射，则映射 \(\mathsf{X}'(h)\colon\mathsf{X}'(\mathsf{B})\to\mathsf{X}'(\mathsf{A})\) 连续。若 h 是满射，则 \(\mathsf{X}'(h)\) 是从 \(\mathsf{X}'(\mathsf{B})\) 到其像上的一个同胚，且这个像在 \(\mathsf{X}'(\mathsf{A})\) 中是闭的。

设 \(A = A_{1} \times \cdots \times A_{n}\)；沿用上面的记号，\(\mathsf{X}'(\pi_{i})\) 是从 \(\mathsf{X}'(\mathrm{A}_{i})\) 到 \(X_{i}'\) 上的一个同胚，且 \(X_{i}'\) 在 \(\mathsf{X}'(\mathrm{A})\) 中是闭的。因此 \(X_{i}' - \{0\}\) 在 \(\mathsf{X}'(\mathrm{A})\) 中是开的；诸 \(\mathsf{X}'(\pi_{i})\) 定义了从诸 \(\mathsf{X}'(\mathrm{A}_{i})\) 的和空间 S 到 \(\mathsf{X}'(\mathrm{A})\) 上的一个连续映射，并且可以立即验证，点 \(0 \in \mathsf{X}'(\mathrm{A}_{1}), \ldots, 0 \in \mathsf{X}'(\mathrm{A}_{n})\) 的邻域之并的像是 \(0 \in \mathsf{X}'(\mathrm{A})\) 的一个邻域；由此得出，\(\mathsf{X}'(\mathrm{A})\) 典范地等同于 S 的一个商空间。特别地，空间 \(\mathsf{X}(\mathrm{A})\) 与诸 \(\mathsf{X}(\mathrm{A}_{i})\) 的和空间视为同一。

从 \(\mathsf{X}^{\prime}(\mathsf{A})\) 到 \(\mathsf{X}(\tilde{\mathsf{A}})\) 的典范双射是一个同胚。设 B 是 K 上的一个幺代数，\(B^{\prime}\) 是其基础代数；则空间 \(\mathsf{X}(\mathsf{B})\) 与子空间 \(\mathsf{X}(\mathsf{B}^{\prime})\)（\(\mathsf{X}^{\prime}(\mathsf{B}^{\prime})\) 的子空间）视为同一。

